% GENERATED FILE. Edit canonical lessons, then run npm run generate:books.
% canonical-source-hash: fnv1a64:0f2b797022805a42
% canonical-lessons: PA-C192-lok-return, PA-C192-lokaan, PA-C192-lokaan-review, PA-C192-lokaan-naal, PA-C192-lokaan-lai, PA-C192-lokaan-ton, PA-C192-payoff

\chapter{People Before Postpositions}
\label{ch:pa-a2-lokaan-oblique}
\hlchaptermodality{192}

\begin{chapteropening}
\textbf{By the end of this chapter:} \emph{I can distinguish the standalone word for people from its form before three familiar postpositions.}

Contrast standalone people with three familiar postposition phrases through listening, speaking, reading, and writing.
\end{chapteropening}

\section[\texorpdfstring{\pa{ਲੋਕ} · \pa{ਨਾਲ} · \pa{ਲਈ} · \pa{ਤੋਂ} (lok; naal; laii; ton)}{ਲੋਕ · ਨਾਲ · ਲਈ · ਤੋਂ (lok; naal; laii; ton)}]{Return to people and three postpositions}

\label{lesson:PA-C192-lok-return}



\begin{quote}
\emph{Hear:} \emph{lok}; YOU SAY: \textbf{\pa{ਲੋਕ}}
\end{quote}

\begin{grammarlens}[title={Grammar Lens}]
Keep the familiar word for “people” ready. Then retrieve the three familiar postpositions \textbf{\pa{ਨਾਲ}}, \textbf{\pa{ਲਈ}}, and \textbf{\pa{ਤੋਂ}} without joining them to the noun yet.
\end{grammarlens}

\subsection*{Your turn}
\begin{itemize}
\raggedright
  \item \emph{Hear:} \emph{lok}; YOU SAY: \textbf{\pa{ਲੋਕ}}
  \item \emph{Hear:} \emph{naal; laii; ton}; YOU SAY: \textbf{\pa{ਨਾਲ}; \pa{ਲਈ}; \pa{ਤੋਂ}}
\end{itemize}

\begin{tcolorbox}[breakable,colback=teal!4,colframe=teal!35!black,title={Before you move on}]
Say “people,” then “with,” “for,” and “from.” This opener owns no new form.
\end{tcolorbox}
\section[\texorpdfstring{\pa{ਲੋਕਾਂ} (lokaan)}{ਲੋਕਾਂ (lokaan)}]{Change one familiar noun before a postposition}

\label{lesson:PA-C192-lokaan}



\begin{quote}
Retrieve \textbf{\pa{ਲੋਕ}}, then the familiar companion postposition \textbf{\pa{ਨਾਲ}}.
\end{quote}

\begin{grammarlens}[title={Grammar Lens: direct owner}]
\begin{quote}
\textbf{\pa{ਲੋਕਾਂ} \pa{ਨਾਲ}} — \emph{lokā̃ nāl} — with people
\end{quote}

The standalone form is \textbf{\pa{ਲੋਕ}}. Before a postposition, this noun becomes \textbf{\pa{ਲੋਕਾਂ}}. Keep the step narrow: this lesson follows one noun and does not say that every Punjabi noun changes in the same way.
\end{grammarlens}

\begin{grammarlens}[title={Grammar Lens: source boundary}]
Johal explains the oblique case before a postposition and shows that this class of masculine plural adds the nasal ending \textbf{-\pa{ਾਂ}}. Here that source rule is applied only to \textbf{\pa{ਲੋਕ}}. Source: Ranbir Johal, \emph{Speaking and Writing Punjabi} (2018), pp. 76–78, \href{https://openlibrary-repo.ecampusontario.ca/jspui/bitstream/123456789/492/2/2018Pun1100CoursePackage.pdf}{oblique nouns before postpositions}.
\end{grammarlens}

\subsection*{Your turn}
\begin{itemize}
\raggedright
  \item \emph{Hear:} \emph{lokaan naal}; YOU SAY: \textbf{\pa{ਲੋਕਾਂ} \pa{ਨਾਲ}}
  \item \emph{Read it:} \textbf{\pa{ਲੋਕਾਂ} \pa{ਨਾਲ}}
  \item \emph{Write it:} copy \textbf{\pa{ਲੋਕਾਂ}} once, then check the final \textbf{\pa{ਾਂ}}
\end{itemize}

\begin{tcolorbox}[breakable,colback=teal!4,colframe=teal!35!black,title={Before you move on}]
Produce the form of “people” used before \textbf{\pa{ਨਾਲ}}.
\end{tcolorbox}
\section[\texorpdfstring{\pa{ਲੋਕਾਂ} (lokaan)}{ਲੋਕਾਂ (lokaan)}]{Retrieve the pre-postposition form}

\label{lesson:PA-C192-lokaan-review}



\begin{quote}
Say the familiar standalone word \textbf{\pa{ਲੋਕ}}.
\end{quote}

\begin{grammarlens}[title={Grammar Lens}]
Keep the contrast small: standalone \textbf{\pa{ਲੋਕ}}; before a postposition \textbf{\pa{ਲੋਕਾਂ}}. No other noun pattern enters this retrieval.
\end{grammarlens}

\subsection*{Your turn}
\begin{itemize}
\raggedright
  \item \emph{Hear:} \emph{lok; lokaan}; YOU SAY: \textbf{\pa{ਲੋਕ}; \pa{ਲੋਕਾਂ}}
  \item \emph{Read it:} \textbf{\pa{ਲੋਕਾਂ}}
  \item \emph{Write it:} cover the model and write \textbf{\pa{ਲੋਕਾਂ}}
\end{itemize}

\begin{tcolorbox}[breakable,colback=teal!4,colframe=teal!35!black,title={Before you move on}]
Produce the form that waits for a postposition.
\end{tcolorbox}
\section[\texorpdfstring{\pa{ਲੋਕਾਂ} \pa{ਨਾਲ} (lokaan naal)}{ਲੋਕਾਂ ਨਾਲ (lokaan naal)}]{Return to the companion frame}

\label{lesson:PA-C192-lokaan-naal}



\begin{quote}
Retrieve the pre-postposition form, then the word for “with.”
\end{quote}

\begin{grammarlens}[title={Grammar Lens: one familiar postposition}]
\begin{quote}
\textbf{\pa{ਲੋਕਾਂ} \pa{ਨਾਲ}} — \emph{lokā̃ nāl} — with people
\end{quote}

The noun comes first. The familiar postposition follows it.
\end{grammarlens}

\subsection*{Your turn}
\begin{itemize}
\raggedright
  \item \emph{Hear:} \emph{lokaan naal}; YOU SAY: \textbf{\pa{ਲੋਕਾਂ} \pa{ਨਾਲ}}
  \item \emph{Read it:} \textbf{\pa{ਲੋਕਾਂ} \pa{ਨਾਲ}}
  \item \emph{Write it:} cover and write \textbf{\pa{ਲੋਕਾਂ} \pa{ਨਾਲ}}
\end{itemize}

\begin{tcolorbox}[breakable,colback=teal!4,colframe=teal!35!black,title={Before you move on}]
Produce “with people.”
\end{tcolorbox}
\section[\texorpdfstring{\pa{ਲੋਕਾਂ} \pa{ਲਈ} (lokaan laii)}{ਲੋਕਾਂ ਲਈ (lokaan laii)}]{Move the same noun form to for}

\label{lesson:PA-C192-lokaan-lai}



\begin{quote}
Retrieve the noun form, then familiar \textbf{\pa{ਲਈ}}.
\end{quote}

\begin{grammarlens}[title={Grammar Lens: keep the noun form}]
\begin{quote}
\textbf{\pa{ਲੋਕਾਂ} \pa{ਲਈ}} — \emph{lokā̃ laī} — for people
\end{quote}

Only the familiar postposition changes. The noun remains \textbf{\pa{ਲੋਕਾਂ}}.
\end{grammarlens}

\subsection*{Your turn}
\begin{itemize}
\raggedright
  \item \emph{Hear:} \emph{lokaan laii}; YOU SAY: \textbf{\pa{ਲੋਕਾਂ} \pa{ਲਈ}}
  \item \emph{Read it:} \textbf{\pa{ਲੋਕਾਂ} \pa{ਲਈ}}
  \item \emph{Write it:} cover and write \textbf{\pa{ਲੋਕਾਂ} \pa{ਲਈ}}
\end{itemize}

\begin{tcolorbox}[breakable,colback=teal!4,colframe=teal!35!black,title={Before you move on}]
Produce “for people.”
\end{tcolorbox}
\section[\texorpdfstring{\pa{ਲੋਕਾਂ} \pa{ਤੋਂ} (lokaan ton)}{ਲੋਕਾਂ ਤੋਂ (lokaan ton)}]{Move the same noun form to from}

\label{lesson:PA-C192-lokaan-ton}



\begin{quote}
Retrieve the noun form, then familiar \textbf{\pa{ਤੋਂ}}.
\end{quote}

\begin{grammarlens}[title={Grammar Lens: a third familiar postposition}]
\begin{quote}
\textbf{\pa{ਲੋਕਾਂ} \pa{ਤੋਂ}} — \emph{lokā̃ tõ} — from people
\end{quote}

The postposition changes again. The pre-postposition noun form stays fixed.
\end{grammarlens}

\subsection*{Your turn}
\begin{itemize}
\raggedright
  \item \emph{Hear:} \emph{lokaan ton}; YOU SAY: \textbf{\pa{ਲੋਕਾਂ} \pa{ਤੋਂ}}
  \item \emph{Read it:} \textbf{\pa{ਲੋਕਾਂ} \pa{ਤੋਂ}}
  \item \emph{Write it:} cover and write \textbf{\pa{ਲੋਕਾਂ} \pa{ਤੋਂ}}
\end{itemize}

\begin{tcolorbox}[breakable,colback=teal!4,colframe=teal!35!black,title={Before you move on}]
Produce “from people.”
\end{tcolorbox}
\section[\texorpdfstring{\pa{ਲੋਕ} · \pa{ਲੋਕਾਂ} \pa{ਨਾਲ} … (lok; lokaan …)}{ਲੋਕ · ਲੋਕਾਂ ਨਾਲ … (lok; lokaan …)}]{A bounded people-and-postpositions payoff}

\label{lesson:PA-C192-payoff}



\begin{quote}
Recall standalone \textbf{\pa{ਲੋਕ}}, then the form used before a postposition.
\end{quote}

\subsection*{Your turn: four skills}
\begin{itemize}
\raggedright
  \item \emph{Hear:} \emph{lok; lokaan naal; lokaan laii; lokaan ton}; YOU SAY: \textbf{\pa{ਲੋਕ}; \pa{ਲੋਕਾਂ} \pa{ਨਾਲ}; \pa{ਲੋਕਾਂ} \pa{ਲਈ}; \pa{ਲੋਕਾਂ} \pa{ਤੋਂ}}
  \item \emph{Read it:} \textbf{\pa{ਲੋਕ} · \pa{ਲੋਕਾਂ} \pa{ਨਾਲ} · \pa{ਲੋਕਾਂ} \pa{ਲਈ} · \pa{ਲੋਕਾਂ} \pa{ਤੋਂ}}
  \item \emph{Write it:} cover the line and write \textbf{\pa{ਲੋਕਾਂ}} once, then one complete postposition phrase
\end{itemize}

\begin{tcolorbox}[breakable,colback=teal!4,colframe=teal!35!black,title={Before you move on}]
Produce the standalone form, then all three postposition phrases.

This is the fifteenth prerequisite tranche exposed by Punjabi A2 mock drafting. It does not unblock Punjabi A2 mock 1 by itself. It is not an A2 mock, not A2 readiness or attainment, and not external or human validation.
\end{tcolorbox}
